\documentclass[a4paper,10pt]{article}
\usepackage[utf8]{inputenc}
\usepackage[T1]{fontenc}
\usepackage{lmodern}
\usepackage[margin=1.4cm, top=1.2cm, bottom=1.2cm]{geometry}
\usepackage{hyperref}
\usepackage{xcolor}
\usepackage{enumitem}
\usepackage{graphicx}
\usepackage{titlesec}

% Minimal clean styling
\hypersetup{colorlinks=true, urlcolor=blue!70!black, linkcolor=blue!70!black}
\pagestyle{empty}
\setlength{\parindent}{0pt}

% Section formatting
\titleformat{\section}{\large\bfseries\color{black}}{}{0em}{}[\titlerule]
\titlespacing*{\section}{0pt}{6pt}{4pt}

% Colors
\definecolor{accent}{RGB}{20, 50, 120}

% Role command
\newcommand{\role}[4]{%
  \vspace{2pt}\noindent
  \begin{minipage}[t]{0.70\textwidth}\raggedright\textbf{#1} \textbullet\ \textit{#3}\end{minipage}%
  \hfill
  \begin{minipage}[t]{0.28\textwidth}\raggedleft\footnotesize\color{gray}#2 \textbullet\ #4\end{minipage}%
  \vspace{1pt}%
}
\setlist[itemize]{leftmargin=1.2em, itemsep=1.5pt, topsep=2pt, parsep=0pt}

\begin{document}

% ================= HEADER =================
\begin{minipage}[c]{0.80\textwidth}
  {\Huge\bfseries Prabhavit Gupta}\\[3pt]
  {\color{accent}\large AI Engineer $\vert$ Production AI Systems, Agent Architectures, Retrieval $\&$ Financial Evals}\\[6pt]
  {\footnotesize +39 351 610 7807 ~$\vert$~
  \href{mailto:prabhavitg@gmail.com}{prabhavitg@gmail.com} ~$\vert$~
  \href{https://www.linkedin.com/in/prabhavit/}{linkedin.com/in/prabhavit} ~$\vert$~
  \href{https://github.com/prab-gupta}{github.com/prab-gupta} ~$\vert$~
  Venezia, Italia (Residente in Italia $\vert$ Permesso di Soggiorno UE a Tempo Indeterminato $\vert$ Full Remote)}
\end{minipage}\hfill

\vspace{6pt}

% ================= EXECUTIVE SUMMARY =================
\section{Profilo Professionale}
\begin{itemize}
  \item \textbf{Sviluppo di Sistemi AI in Produzione:} Software engineer con ~3 anni di esperienza commerciale nello sviluppo e rilascio in produzione di sistemi basati su LLM, architetture agentiche e pipeline di retrieval per flussi di dati finanziari e aziendali. Forte focus su affidabilità deterministica, automazione end-to-end e codice pronto per la produzione (FastAPI, Docker, CI/CD, no prototipi usa-e-getta).
  \item \textbf{Knowledge Management, Re-ranking e Tool Use:} Ampia esperienza nella progettazione di sistemi di semantic search, vector databases (pgvector), re-ranking cross-encoder e integrazione di tool sicuri tramite Model Context Protocol (MCP Server \& Client). Specializzato nell'interrogazione strutturata di dati transazionali, fatture e contabilità.
  \item \textbf{Framework di Valutazione Rigorosi (Evals):} Padroneggia la definizione e il monitoraggio di test suite di valutazione offline e online (RAGAS, test di regressione automatizzati, Langfuse tracing), quantificando sistematicamente accuratezza, aderenza ai fatti, latenza d'inferenza e costi computazionali prima di ogni deploy.
\end{itemize}

% ================= TECHNICAL SKILLS =================
\section{Competenze Tecniche}
\begin{itemize}[leftmargin=1.2em, itemsep=2pt]
  \item \textbf{AI \& LLM Engineering:} Production LLM Systems, Agentic Workflows, Tool Calling, Knowledge Retrieval \& Memory, Semantic Search \& Re-ranking, LangGraph, Prompt Engineering, Structured Outputs (Pydantic).
  \item \textbf{Data \& Protocols:} Model Context Protocol (MCP Server \& Client SDKs), PostgreSQL (pgvector), Redis (Caching), Data Pipelines, Document Parsing, Financial Transaction Classification.
  \item \textbf{Evals \& Observability:} Eval-Driven Development, RAGAS Benchmarking, Regression Testing, Langfuse Distributed Tracing, Latency \& Token Cost Profiling, Continuous Monitoring.
  \item \textbf{Software Engineering:} Python moderno (AsyncIO, FastAPI, PyTest, Typing), TypeScript, Elixir basics, Docker, Microservices, REST \& WebSocket APIs, Git, GitHub Actions CI/CD.
  \item \textbf{Metodologie:} High Ownership, Architettura Difensiva, Code Reviews, ADRs, Collaborazione con Product, Data ed Engineering teams.
\end{itemize}

% ================= PROFESSIONAL EXPERIENCE =================
\section{Esperienza Professionale}

\role{AI Engineer \& Consulente Tecnico}{2024 -- Presente}{Consulenza Indipendente di Ingegneria AI}{Venezia, IT (Remoto)}
\begin{itemize}
  \item Progettato e rilasciato in produzione sistemi AI agentici e motori di retrieval aumentato per applicazioni enterprise europee, curando l'intero ciclo di vita dall'analisi dei requisiti al deploy.
  \item Implementato pipeline avanzate di retrieval, context management e re-ranking su banche dati relazionali e vettoriali, migliorando la precisione di estrazione e riconciliazione su documenti amministrativi e transazioni.
  \item Sviluppato server e client custom conformi al Model Context Protocol (MCP), consentendo ad agenti autonomi di eseguire tool use sicuro su database PostgreSQL con controllo rigoroso dei permessi e audit trail.
  \item Istituito framework sistematici di evaluation (evals) con dataset di test dedicati e benchmark RAGAS per validare metriche di completezza, precisione e controllo delle allucinazioni prima del rilascio.
  \item Integrato tracciamento distribuito (Langfuse, OpenTelemetry) per monitorare latenze p95/p99, consumi di token e comportamenti anomali su pipeline in produzione.
  \item Containerizzato microservizi in Docker e automatizzato i flussi di continuous integration e testing regressionale via GitHub Actions.
\end{itemize}

\vspace{3pt}
\role{Full-Stack AI Software Engineer}{2023 -- 2024}{Tech Automation Systems}{Milano / Remoto, IT}
\begin{itemize}
  \item Sviluppato servizi backend asincroni in Python e TypeScript per l'estrazione, la normalizzazione e la classificazione automatica di elevati volumi di dati aziendali.
  \item Implementato pipeline di ricerca semantica con PostgreSQL e Redis, garantendo tempi di risposta inferiori al secondo e caching efficiente degli embeddings.
  \item Scritto suite di test esaustive con PyTest e introdotto verifiche automatizzate di qualità del codice nei processi di integrazione continua.
  \item Collaborato attivamente con product manager, designer e stakeholder di business per tradurre requisiti operativi in soluzioni software robuste e modulari.
\end{itemize}

% ================= NOTABLE PROJECTS =================
\section{Progetti Rilevanti}
\begin{itemize}
  \item \textbf{Autonomous Tool-Using Agent Platform (Python + LangGraph + MCP + Postgres):} Piattaforma agentica multi-fase coordinata con gestione dello stato, memoria condivisa, integrazione tool via MCP e tracciamento distribuito Langfuse.
  \item \textbf{Motore di Analisi Dati e Documenti Finanziari (FastAPI + pgvector):} Microservizio di retrieval semantico e parsing strutturato con convalida Pydantic, algoritmi di re-ranking e harness di valutazione per eliminare le regressioni.
\end{itemize}

% ================= EDUCATION & LEGAL STATUS =================
\section{Formazione, Lingue e Stato Legale}
\begin{itemize}
  \item \textbf{Laurea Triennale in Economia e Filosofia (Philosophy \& Economics)} -- Università Ca' Foscari di Venezia (2018 -- 2022) \\
  \textit{Solida formazione quantitativa e teorica in logica formale, modelli decisionali, statistica e analisi dei sistemi economici complessi.}
  \item \textbf{Lingue:} Italiano (Fluente / Professionale), Inglese (C1 Avanzato / Full Professional).
  \item \textbf{Stato Legale:} Permesso di Soggiorno UE per Soggiornanti di Lungo Periodo a Tempo Indeterminato (Italia) -- Pieno diritto di lavoro e residenza permanente in Italia senza necessità di visto o sponsorizzazione; disponibile immediatamente.
\end{itemize}

\end{document}
